\documentclass[../../../include/open-logic-section]{subfiles}

\begin{document}

\olfileid{sth}{z}{pairs}
\olsection{અક્રમયુક્ત જોડ}

આગળ વિચારવાની સ્વયંસિદ્ધિ આ છે:

\begin{axiom}[અક્રમયુક્ત જોડ]
કોઈ પણ ગણો $a, b$ માટે ગણ $\{a, b\}$ અસ્તિત્વ ધરાવે છે.
\[
\forall a \forall b \exists P \forall x (x \in P \liff (x = a \lor x = b))
\]
\end{axiom}

પુનરાવર્તી કલ્પના વડે આ સ્વયંસિદ્ધિને આ રીતે યોગ્ય ઠરાવી શકીએ.
ધારો કે $a$ તબક્કા $S$ પર અને $b$ તબક્કા $T$ પર ઉપલબ્ધ છે.
તબક્કાઓ $S$ અને $T$માંથી જે પાછળ આવે તેને $M$ લો. તો $a$ અને
$b$ બંને તબક્કા $M$ પર ઉપલબ્ધ હોવાથી $\{a,b\}$ એવો સંગ્રહ
$M$ પછીના કોઈ પણ તબક્કે રચી શકાય (આ બે તબક્કાઓમાંથી જે પાછળ હોય
તે લીધેલો છે).

પરંતુ થોભો!{} $M$ પછી કોઈ તબક્કાઓ \emph{છે} એવું શા માટે માનવું?
જો કોઈ ન હોય, તો આપણું ઔચિત્ય નિષ્ફળ જશે. તેથી અક્રમયુક્ત જોડની
સ્વયંસિદ્ધિને યોગ્ય ઠરાવવા માટે \olref[sth][z][story]{sec}માં
આપેલા વર્ણનમાં વધુ એક સિદ્ધાંત ઉમેરવો પડે:
\begin{enumerate}
\item[] \stagessucc. કોઈ અંતિમ તબક્કો નથી.
\end{enumerate}
શું આ સિદ્ધાંત યોગ્ય છે? શોએનફેલ્ડના વર્ણનમાં કોઈ અંતિમ તબક્કો
નથી એવું \emph{સ્પષ્ટ રીતે} ક્યાંય કહ્યું નહોતું. છતાં,
ચોક્કસ રીતે કહીએ તો આ આપણા વર્ણનમાં વધારાનો ઉમેરો હોય, તો પણ
ગણો તબક્કાઓમાં રચાય છે તે મૂળભૂત વિચાર સાથે તે સારી રીતે બંધબેસે
છે. આગળ આપણે તેને સ્વીકારી લઈશું. તેથી અક્રમયુક્ત જોડની
સ્વયંસિદ્ધિ પણ સ્વીકારીશું.

આ નવી સ્વયંસિદ્ધિ મળતાં બીજા ઘણા ગણોનું અસ્તિત્વ સાબિત કરી
શકીએ છીએ. દાખલા તરીકે:

\begin{prop}\ollabel{prop:pairsconsequences}
કોઈ પણ ગણો $a$ અને $b$ માટે નીચેના ગણો અસ્તિત્વ ધરાવે છે:
\begin{enumerate}
\item\ollabel{singleton} $\{a\}$
\item\ollabel{binunion} $a \cup b$
\item\ollabel{tuples} $\tuple{a, b}$
\end{enumerate}
\end{prop}

\begin{proof}
\olref{singleton}. અક્રમયુક્ત જોડની સ્વયંસિદ્ધિથી $\{a, a\}$
અસ્તિત્વ ધરાવે છે, જે ઘટકો દ્વારા નિર્ધારિત સમાનતાને કારણે
$\{a\}$ જ છે.

\olref{binunion}. અક્રમયુક્ત જોડની સ્વયંસિદ્ધિથી $\{a, b\}$
અસ્તિત્વ ધરાવે છે. હવે યોગગણની સ્વયંસિદ્ધિથી
$a \cup b = \bigcup
\{a, b\}$ અસ્તિત્વ ધરાવે છે.

\olref{tuples}. \olref{singleton}થી $\{a\}$ અસ્તિત્વ ધરાવે છે.
અક્રમયુક્ત જોડની સ્વયંસિદ્ધિથી $\{a,
b\}$ અસ્તિત્વ ધરાવે છે. હવે ફરીથી એ જ સ્વયંસિદ્ધિથી
$\{\{a\}, \{a, b\}\} = \tuple{a, b}$ અસ્તિત્વ ધરાવે છે.
\end{proof}

\begin{prob}
કોઈ પણ ગણો $a, b, c$ માટે ગણ $\{a, b, c\}$ અસ્તિત્વ ધરાવે છે
એ બતાવો.
\end{prob}

\begin{prob}
કોઈ પણ ગણો $a_1, \ldots, a_n$ માટે ગણ $\{a_1, \ldots,
a_n\}$ અસ્તિત્વ ધરાવે છે એ બતાવો.
\end{prob}

%\begin{proof} By two applications of Pairs, $\{\{a_1, a_2\}, \{a_1,
%   a_3\}\}$ exists. By Union and Extensionality, $\bigcup \{\{a_1,
%   a_2\}, \{a_1, a_3\}\} = \{a_1, a_2, a_3\}$ exists. Repeat this
%   trick as often as necessary. \end{proof}

\end{document}
